%% ============================================================
%%  Introduction générale
%% ============================================================
\chapter*{Introduction générale}
\addcontentsline{toc}{chapter}{Introduction générale}
\markboth{Introduction générale}{Introduction générale}

\section*{Contexte et motivation}

Le trading sur les marchés financiers est une activité qui se joue en quelques secondes.
Sur la Salle des Marchés de \textbf{BMCE Capital}, l'un des acteurs majeurs du marché
des capitaux marocain, les équipes d'Equity Trading sont confrontées quotidiennement à
un volume d'information considérable~: graphiques prix, indicateurs techniques, données
de volume, actualités macroéconomiques, résultats d'entreprises. Dans ce contexte,
la capacité à \emph{lire le marché rapidement, justifier une décision et définir des
niveaux d'entrée et de sortie précis} constitue un avantage compétitif direct.

\medskip

Parallèlement, la quantification des stratégies de trading — c'est-à-dire leur
formalisation mathématique et leur validation sur données historiques — est devenue
une pratique incontournable dans les institutions financières de premier rang. Pourtant,
cette validation reste piégée par un problème fondamental bien documenté dans la
littérature scientifique~: l'overfitting, ou sur-ajustement. Une stratégie qui
fonctionne parfaitement sur des données historiques peut se révéler totalement
inefficace en production, non pas parce qu'elle est mal conçue, mais parce qu'elle
a été \og calibrée \fg{} sur ces mêmes données.

\section*{Problématique}

Ce projet de fin d'études répond à une double problématique, intimement liée aux
enjeux opérationnels de BMCE Capital~:

\begin{enumerate}
    \item \textbf{Comment identifier les meilleures opportunités de marché} parmi les
    93 titres cotés au MASI, avec suffisamment de rapidité et de traçabilité pour être
    actionnable en conditions réelles de trading~?

    \item \textbf{Comment valider rigoureusement des stratégies de trading}
    sans tomber dans le piège de l'overfitting, du data-snooping ou du biais de
    sélection sur tests multiples~?
\end{enumerate}

Ces deux questions ne peuvent pas être traitées indépendamment~: une opportunité de
marché n'est pertinente que si elle repose sur des signaux dont la robustesse a été
démontrée hors-échantillon, et une stratégie valide ne génère de valeur que si
elle est lisible et actiona\-ble par le trader.

\section*{Objectifs}

Le projet vise à concevoir, développer et déployer une \textbf{plateforme quantitative
complète} structurée en quatre stations fonctionnelles~:

\begin{enumerate}
    \item \textbf{Données}~: gestion de l'univers de marché (OHLCV, 93 tickers MASI,
    fraîcheur automatique).
    \item \textbf{Signaux}~: découverte et évaluation OOS des indicateurs techniques
    via un pipeline à 7 couches.
    \item \textbf{Stratégie}~: définition des règles de trading et paramétrage
    pour l'optimisation.
    \item \textbf{Backtest}~: validation par Walk-Forward Analysis avec critère
    d'acceptation WFE\,$\geq 50\,\%$.
\end{enumerate}

Le livrable final est un \textbf{tableau de bord décisionnel} qui consomme les deux
moteurs de résultat (signal engine + WFO) pour fournir, par titre, une recommandation
directionnelle avec niveaux et conviction.

\section*{Méthodologie}

L'approche méthodologique s'inscrit dans le cadre du paradigme méta-stratégie de
López de Prado~\cite{lopezdeprado2018} et de la Walk-Forward Analysis de Pardo~\cite{pardo2008}.
Elle est structurée autour de trois principes~:

\begin{itemize}
    \item \textbf{Séparation stricte des stations}~: données, signaux, stratégie
    et backtest sont des modules indépendants pour prévenir l'overfitting cognitif.
    \item \textbf{Évaluation systématiquement hors-échantillon}~: aucun résultat
    in-sample n'est présenté comme preuve de performance.
    \item \textbf{Multi-défenses anti-overfitting}~: chaque couche du pipeline
    intègre un mécanisme de contrôle spécifique (gate de viabilité, score
    multi-composante, clustering de corrélation, WFE).
\end{itemize}

\section*{Organisation du mémoire}

Ce rapport est organisé en sept chapitres~:

\begin{itemize}
    \item Le \textbf{chapitre~1} pose le contexte industriel (BMCE Capital, marché
    MASI), formule la double problématique et établit l'état de l'art des solutions
    existantes.

    \item Le \textbf{chapitre~2} présente les fondements de l'analyse technique~:
    les vingt familles d'indicateurs (Tendance, Momentum, Oscillation, Volume),
    leurs formules mathématiques et leurs limites structurelles qui motivent
    les chapitres suivants.

    \item Le \textbf{chapitre~3} détaille le premier moteur de résultat~: le
    Moteur de Signaux Personnalisé, pipeline de sept couches (A→G) évaluant
    20 familles d'indicateurs par validation de stabilité temporelle sur sous-périodes.

    \item Le \textbf{chapitre~4} présente le second moteur~: le Moteur WFO des
    Signaux, qui applique la Walk-Forward Analysis de Pardo à la sélection dynamique
    d'indicateurs par catégorie, avec la fonction objectif PROM et le système de
    notation A–F.

    \item Le \textbf{chapitre~5} expose l'architecture technique de la plateforme~:
    stack full-stack, quatre stations déployées (Données, Signaux, Dashboard,
    Analytics) et bonnes pratiques DevOps.

    \item Le \textbf{chapitre~6} présente le livrable décisionnel~: le tableau de
    bord unifié, sa réponse à la problématique, les résultats expérimentaux sur
    données MASI et l'évaluation opérationnelle.

    \item Le \textbf{chapitre~7} apporte la preuve statistique de l'edge~: matrice
    de capacité prédictive, courbes IC, Deflated Sharpe Ratio et contrôle du taux
    de fausses découvertes (FDR).
\end{itemize}

\vspace{0.5cm}
\noindent La conclusion synthétise les contributions, identifie les limites et
trace les perspectives de développement.
